\chapter{\objectivesname} \label{chp:objetivos}

En este capítulo se detalla el objetivo principal planteado para el desarrollo de este trabajo, así como los objetivos específicos en los que se desglosa para poder alcanzar dicho propósito. Además, se presentan las contribuciones más relevantes que se han aportado tras la realización de este Trabajo de Fin de Grado (TFG).

\section{Objetivos}

El principal objetivo de este trabajo es proponer una solución que permita reducir las emisiones de $CO_2$ en redes cableadas provinientes del consumo energético de los dispositivos de red. Haciendo uso de modelos de enrutamiento del tráfico de datos y mediante el diseño de un algoritmo de apagado de enlaces, se modifica la configuración de la red para reducir el consumo energético de los dispositivos de red y, por tanto, las emisiones de $CO_2$ asociadas a dicho consumo. Para ello, se han definido los siguientes objetivos específicos:

\begin{enumerate}
    \item Análisis de datos y modelado de la red.
    \begin{itemize}
        \item Identificar y procesar las fuentes de datos necesarias, como la API de SNDlib \cite{Orlowski2010}, para modelar las redes cableadas y sus emisiones de $CO_2$.
        \item Representar la topología de red mediante estructuras de datos adecuadas, como matrices o listas, con las que trabajar a partir de archivos \textit{csv} y \textit{json}.
    \end{itemize}
    \item Estudio del apagado de enlaces.
    \begin{itemize}
        \item Investigar el funcionamiento de la técnica de apagado de enlaces para la reducción del consumo energético en redes SDN.
        \item Estudiar criterios y restricciones a tener en cuenta para el apagado de enlaces sin comprometer la conectividad o el rendimiento de la red.
        \item Analizar el impacto del apagado de enlaces sobre el enrutamiento del tráfico de datos y las emisiones de $CO_2$.
    \end{itemize}
    \item Estudio de soluciones de enrutado conscientes de carbono.
    \begin{itemize}
        \item Estudiar diversas técnicas de encaminamiento del tráfico de datos.
        \item Analizar diferentes propuestas de enrutamiento conscientes de carbono.
        \item Estudiar cómo incorporar la información de las emisiones en el proceso de enrutamiento.
    \end{itemize}
    \item Desarrollo del algoritmo de optimización.
    \begin{itemize}
        \item Diseñar e implementar un algoritmo de PSO que se adapte a las características del problema planteado.
        \item Integrar técnicas de modularización y reutilización de código para facilitar el desarrollo y la escalabilidad del algoritmo de cara a futuras mejoras.
    \end{itemize}
    \item Validación y análisis de resultados.
    \begin{itemize}
        \item Hacer uso de topologías reales para evaluar el rendimiento del algoritmo desarrollado.
        \item Analizar los resultados obtenidos para determinar la efectividad del algoritmo en la reducción de emisiones de carbono.
    \end{itemize}
\end{enumerate}

\section{Contribuciones}

Tras la realización de este trabajo, se han aportado las contribuciones que a continuación se detallan:

\begin{itemize}
    \item 
\end{itemize}